\paragraph{Learning problem.}
Let $q$ denote a question and $X$ denote a patient record represented as a set of candidate chunks $c \in X$ (note chunks and event snippets). Our goal is to learn a scoring function $s_\theta(q,c)$ such that selecting the top-$K$ chunks by score yields a compact context that enables a frozen downstream LLM to answer $q$ accurately.

\paragraph{Chunk scoring model.}
We train an L1-regularized logistic regression model to predict chunk usefulness:
\[
p_\theta(y=1 \mid q,c) = \sigma\left(w^\top \phi(q,c)\right),
\]
where $\phi(q,c)$ is the engineered feature vector and $\sigma(\cdot)$ is the sigmoid. L1 regularization induces sparsity in $w$, improving interpretability and effectively performing feature selection over the interaction features .

\paragraph{Training labels and objective.}
We define $y=1$ if the chunk is deemed answer-supporting under an overlap heuristic (token-F1 between the chunk text and the answer text exceeds a threshold), and $y=0$ otherwise . The model is trained by minimizing the regularized negative log-likelihood on (question, chunk) examples, with class-weighting and negative sampling to address class imbalance.

\paragraph{Train/validation protocol.}
To avoid leakage from patient-specific wording, we split at the patient level: all examples from a patient are assigned entirely to train or validation . TF--IDF is fit on train-only text and applied to validation via \texttt{transform}; embeddings are computed with a frozen encoder .

\paragraph{Inference-time selection.}
Given a new $(q,X)$ pair, we compute $p_\theta(y=1\mid q,c)$ for all chunks $c \in X$, rank chunks, and select the top-$K$ (or a token-budgeted set) to pass to the downstream LLM. We also consider coherence post-processing (e.g., prefer contiguous windows in notes or coherent time windows in events) as an optional ablation.

\paragraph{Downstream evaluation framework.}
To isolate retrieval quality, we keep the downstream LLM and prompt format fixed and compare retrieval strategies that differ only in context selection. The evaluation pipeline implements six strategies (discharge-only, full-context, recency, BM25, semantic RAG, learned selector), uses a response cache for reproducibility, and supports both OpenAI and HuggingFace models. Importantly, we use a hold out dataset, the EHRQA dataset of 962 clinicial reviewed QA pairs, as evaluation \cite{PhysioNet-ehr-notes-qa-llms-1.0.1}. An important distinction is that while we train on QA pairs generated by a larger model, we evaluate on QA pairs from real physicians.